\section{栈设置与 \texttt{kmain} 调用}

进入高半区后，第一件事是设置栈和调用 C 入口：

\codefile{src/boot/boot.asm}
\begin{lstlisting}[style=nasmstyle]
section .text
_start_high:
    mov esp, stack_top      ; 栈顶（高半区虚拟地址）
    and esp, 0xFFFFFFF0     ; 16 字节对齐（SSE 要求）

    ; cdecl 调用约定：参数从右向左压栈
    push edi                ; mb2_info (物理地址)
    push esi                ; mb2_magic
    call kmain
    add esp, 8              ; 清理参数（不会执行到这里）

.hlt:
    hlt
    jmp .hlt

section .bss
align 4096
stack_bottom:
    resb 16384              ; 16 KiB 内核栈
stack_top:
\end{lstlisting}

\begin{itemize}
  \item 栈在 \texttt{.bss} 段中分配 16\,KiB（\texttt{resb 16384}），
        \texttt{stack\_top} 是栈的高地址端（x86 栈向下增长）。
  \item \texttt{and esp, 0xFFFFFFF0} 确保 16 字节对齐——System V ABI 要求
        调用函数前 \reg{ESP} 必须 16 字节对齐。
  \item \texttt{push edi; push esi; call kmain}——\texttt{cdecl} 约定，
        参数从右向左压栈，\texttt{kmain(uint32\_t magic, void* mb2\_info)}。
  \item \texttt{kmain} 正常不会返回；如果返回则 \texttt{hlt} 停住 CPU。
\end{itemize}

% ============================================================================

